\chapter{Literature Survey}
\label{ch:literature_survey}

\section{Introduction}
The realization of an autonomous, privacy-preserving desktop voice intelligence system lies at the confluence of four major computer science domains:
\begin{enumerate}
    \item Acoustic Signal Processing and Automated Speech Recognition (ASR).
    \item Large Language Model (LLM) Agentic Architectures and Autonomous Code Synthesis.
    \item Operating System Automation, Subprocess Sandboxing, and Computer Use.
    \item Workstation Privacy, Human-in-the-Loop (HITL) Verification, and Access Control.
\end{enumerate}

This chapter conducts a rigorous scholarly survey of foundational and contemporary literature governing these fields. It analyzes the architectural paradigms, mathematical formulations, and engineering limitations of prevailing solutions, establishing the theoretical and practical motivation for PRIVACY68.

\section{Speech Recognition and Acoustic Signal Conditioning}

\subsection{Evolution of Automated Speech Recognition Architectures}
Automated Speech Recognition (ASR) has undergone three fundamental architectural epochs:

\subsubsection{Hidden Markov Models and Gaussian Mixture Models (HMM-GMM)}
Early conversational systems relied on statistical generative models combining Hidden Markov Models for temporal acoustic modeling and Gaussian Mixture Models for acoustic feature representation \cite{schuster2018speech}. These systems required handcrafted acoustic feature representations (e.g., Mel-Frequency Cepstral Coefficients, MFCCs) coupled with static n-gram language models. While computationally lightweight, HMM-GMM pipelines exhibited extreme vulnerability to acoustic reverberation, ambient room noise, and speaker accent variations, achieving Word Error Rates (WER) often exceeding 25\% to 35\% in unconstrained domestic environments.

\subsubsection{Connectionist Temporal Classification (CTC) and Hybrid Neural Models}
The advent of deep learning introduced end-to-end neural acoustic models. Graves et al. \cite{graves2006connectionist} formulated Connectionist Temporal Classification (CTC), an alignment-free loss function that allowed recurrent neural networks (RNNs and LSTMs) to directly map continuous acoustic feature sequences to variable-length character transcriptions without requiring intermediate frame-level phoneme alignments. While CTC significantly improved recognition speed, its conditional independence assumption across output labels limited its ability to model complex long-range linguistic context.

\subsubsection{Attention-Based Encoder-Decoder (AED) and Transformer Architectures}
The Transformer architecture introduced by Vaswani et al. \cite{vaswani2017attention} replaced recurrence with multi-head self-attention mechanisms, enabling parallel acoustic processing and superior long-range semantic representation. Gulati et al. \cite{gulati2020conformer} combined self-attention with depthwise separable convolutions in the Conformer architecture, capturing both global context and local acoustic correlations. Concurrently, Baevski et al. \cite{baevski2020wav2vec} introduced wav2vec 2.0, demonstrating that self-supervised pre-training on thousands of hours of unlabelled speech could yield state-of-the-art representations with minimal fine-tuning.

Radford et al. \cite{radford2023robust} advanced this frontier with Whisper, an attention-based sequence-to-sequence model trained on 680,000 hours of weakly supervised multilingual acoustic data. Whisper demonstrated unprecedented zero-shot robustness across dialects, background noise, and technical jargon. However, deploying full-precision Whisper models locally introduces significant computational bottlenecks. On standard consumer CPUs without high-end discrete GPUs, processing a 3-second audio utterance via Whisper-large can incur 1,500ms to 3,000ms of latency, severely impairing conversational interactivity. While quantized inference engines (such as Faster-Whisper using CTranslate2 int8 quantization) reduce inference latency to 200--400ms, they remain highly susceptible to non-linear microphone gain variations, low Signal-to-Noise Ratios (SNR), and numerical clipping.

\subsection{Acoustic Gain Variability and Peak Amplitude Normalization}
A ubiquitous challenge in consumer desktop environments is the physical heterogeneity of audio capture transducers. Unlike dedicated smart speakers equipped with calibrated circular microphone arrays and hardware DSP chips, desktop workstations connect to diverse audio hardware: low-cost analog 3.5mm headset microphones, integrated laptop webcams, and professional studio condenser microphones.

Desai, Raman, and Narayanan \cite{desai2021speaker} investigated the degradation of neural ASR confidence scores under non-calibrated input conditions. Their findings established that speech captured below $-24\text{ dBFS}$ results in a 41.2\% relative increase in WER due to sub-optimal quantization depth during 16-bit to 32-bit floating point conversion. Conversely, acoustic peaks exceeding $0\text{ dBFS}$ cause severe non-linear harmonic distortion due to analog-to-digital converter (ADC) clipping, completely obliterating high-frequency formant features.

While commercial voice platforms employ adaptive Automatic Gain Control (AGC) algorithms, traditional time-domain AGC introduces continuous dynamic modulation, creating transient ``breathing'' noise artifacts that confuse neural feature extractors. Desai et al. demonstrated that linear peak amplitude normalization:
\begin{equation}
    x_{\text{normalized}}[t] = \left( \frac{x[t]}{\max_{k} |x[k]|} \right) \cdot \alpha
\end{equation}
where $\alpha \in [0.85, 0.95]$ provides optimal crest factor preservation without introducing non-linear dynamic distortion. By maintaining a target headroom of 10\% ($\alpha = 0.90$), linear peak normalization standardizes disparate microphone inputs into an optimal dynamic window for both cloud and local neural acoustic models.

\section{Large Language Models as Autonomous Agents}

\subsection{The Reasoning and Acting (ReAct) Paradigm}
Large Language Models have evolved from statistical sequence predictors \cite{sutskever2014sequence, brown2020language} into versatile general-purpose reasoning engines \cite{bubeck2023sparks}. However, standalone language models remain epistemologically closed: their knowledge is frozen at training time, and they cannot interact with external software or observe operating system states.

To overcome this limitation, Yao et al. \cite{yao2023react} pioneered the \textbf{ReAct} (Reasoning and Acting) framework. ReAct synergizes chain-of-thought reasoning with actionable tool invocations in an iterative cycle:
\begin{equation}
    \text{Thought}_t \longrightarrow \text{Action}_t \longrightarrow \text{Observation}_t \longrightarrow \text{Thought}_{t+1}
\end{equation}
By interleaving reasoning traces with discrete tool executions, ReAct enables models to perform dynamic planning, query APIs, retrieve environment observations, and adjust future actions based on real-time feedback.

Wang et al. \cite{wang2023survey} synthesized the agentic landscape into four core structural pillars:
\begin{enumerate}
    \item \textbf{Profile Module:} Defines the agent's persona, operational boundaries, and system prompt constraints.
    \item \textbf{Memory Module:} Manages short-term conversational context and long-term persistent storage (e.g., vector embeddings or structured databases).
    \item \textbf{Planning Module:} Decomposes complex high-level goals into sequential executable sub-tasks using methods such as Chain-of-Thought (CoT) \cite{wei2022chain} or Tree-of-Thoughts.
    \item \textbf{Action Module:} Translates planned tasks into external environment invocations, such as API calls, web browsing actions, or shell commands.
\end{enumerate}

\subsection{Tool Use and Autonomous Code Synthesis}
Schick et al. \cite{schick2023toolformer} demonstrated in \textbf{Toolformer} that language models can self-supervisedly teach themselves to decide which external APIs to call, when to call them, and how to parse their arguments. In desktop operating system environments, however, pre-defined API catalogs are fundamentally inadequate. An operating system offers thousands of potential interactions---manipulating file systems, querying process trees, adjusting audio mixers, positioning windows, and communicating over network sockets. Constraining an agent to a fixed set of JSON function-calling definitions creates severe inflexibility.

Consequently, modern computer use agents favor \textbf{dynamic code synthesis}. Chen et al. \cite{chen2021evaluating} (OpenAI Codex) and Rozi{\`e}re et al. \cite{roziere2023code} (Code Llama) proved that foundation models trained on extensive source code repositories exhibit extraordinary proficiency in translating natural language directives into syntactically valid, multi-line Python scripts. Writing Python automation code provides maximum expressive power: the agent can leverage standard libraries (\texttt{os}, \texttt{sys}, \texttt{subprocess}, \texttt{shutil}, \texttt{winreg}) and third-party automation packages to orchestrate arbitrary desktop workflows.

\subsection{Self-Healing and Reflective Error Correction}
Despite their coding proficiency, generative models remain prone to syntax hallucinations, obsolete API parameter usage, and runtime type mismatches \cite{geirhos2020beyond}. When executing generated code directly against the host operating system, an unhandled exception or non-zero exit code halts the workflow.

To resolve this fragility, Shinn et al. \cite{shinn2023reflexion} introduced \textbf{Reflexion}, a framework for verbal reinforcement learning where agents reflect on environment execution errors. Rather than updating neural weights, Reflexion maintains an execution traceback memory:
\begin{equation}
    \text{Context}_{t+1} = \text{Context}_t \cup \{\text{Script}_t, \text{Stderr}_t, \text{ExitCode}_t\}
\end{equation}
When the language model receives the raw execution traceback and stderr stream, it identifies the failure locus (e.g., an incorrect library import or unescaped file path) and synthesizes an amended script. In PRIVACY68, this self-healing execution loop operates autonomously, retrying failed scripts up to three iterations before notifying the user.

\subsection{Multi-Provider Model Orchestration and Resilience}
Enterprise and workstation agent deployments frequently suffer from vendor lock-in. Relying on a single proprietary foundation model API exposes the application to severe operational vulnerabilities: sudden rate limits (HTTP 429), regional outages (HTTP 503), latency spikes, and deprecation of specific model versions.

Berquand and Higgins \cite{litellm2024} introduced \textbf{LiteLLM}, an open-source abstraction gateway that unifies over one hundred commercial and open-weights LLM providers under an OpenAI-compatible interface. By standardizing wire formats, token tracking, and authentication headers, LiteLLM enables dynamic model cascading: if a primary high-performance cloud model (e.g., Google Gemini 2.5 Flash) encounters a timeout or rate limit, the request is instantaneously routed to an alternative provider (e.g., Groq Cloud LPU or OpenRouter) or a local offline instance (Ollama). PRIVACY68 leverages LiteLLM as its core reasoning switchboard, providing uninterrupted conversational intelligence.

\section{Workstation Automation and Security Sandboxing}

\subsection{Desktop Computer Use and Benchmarking}
Automating desktop workstations using AI agents represents the modern evolution of Robotic Process Automation (RPA). Traditional RPA platforms (e.g., UiPath, Automation Anywhere) depend on brittle UI selector paths, fixed pixel coordinates, or rigid optical character recognition (OCR) rules that break upon minor GUI redesigns.

Recently, Xie et al. \cite{xie2024osworld} introduced \textbf{OSWorld}, an extensive benchmark evaluating multimodal AI agents across real operating system tasks in Ubuntu and Windows environments. Their empirical findings revealed that while current agents achieve moderate success on single-application tasks, their autonomous execution reliability drops significantly when managing multi-step, multi-application workflows. Furthermore, arbitrary script execution frequently exposes host systems to unrecoverable side effects.

\subsection{Security Threats in Autonomous Code Execution}
Executing LLM-generated code with ambient user privileges introduces critical security threats:
\begin{enumerate}
    \item \textbf{Indirect Prompt Injection:} Carlini et al. \cite{carlini2024prompt} demonstrated that malicious instructions embedded in untrusted external data (e.g., email bodies, web pages, or repository README files) can hijack an agent's reasoning loop, prompting it to exfiltrate private credentials or execute destructive commands.
    \item \textbf{Accidental System Destruction:} Generative models lack intrinsic awareness of physical data persistence. A natural request such as ``clean up my old temporary folders'' can be hallucinated into an unconstrained recursive directory wipe (\texttt{shutil.rmtree('C:\textbackslash')}), causing catastrophic data loss \cite{shieh2020security}.
    \item \textbf{Subprocess Deadlocks and Runaway Processes:} Code generated by an LLM may include infinite loops or blocking socket operations that permanently consume CPU cycles or freeze desktop UI threads.
\end{enumerate}

To protect against these threats, autonomous workstation agents require a multi-tiered security defense:
\begin{itemize}
    \item \textbf{Subprocess Isolation with Strict Timeouts:} Generated automation scripts must execute within dedicated child processes (\texttt{subprocess.run}) governed by hard execution timeouts (e.g., 30 seconds) and piped standard streams.
    \item \textbf{Static Code Inspection:} Pre-execution analysis using Abstract Syntax Trees (AST) or deterministic regex guards to detect forbidden system calls (e.g., format commands, registry modifications, raw socket creation).
    \item \textbf{Human-in-the-Loop (HITL) Tiered Gating:} Halting sensitive or destructive operations behind an interactive, time-bounded user confirmation dialogue \cite{malkin2022privacy}.
\end{itemize}

\section{Privacy Paradigms: Always-On vs. Push-to-Talk}
The foundational ethos of PRIVACY68 is the absolute rejection of the always-on acoustic monitoring paradigm.

Schuster et al. \cite{schuster2018speech} and Rawassizadeh et al. \cite{rawassizadeh2015smart} conducted extensive telemetry analyses of consumer smart speakers, revealing that proprietary voice assistants transmit compressed audio frames to cloud servers even during un-triggered idle periods to calibrate background noise baselines. In corporate environments, law firms, healthcare facilities, and software development workstations, this continuous surveillance model violates basic confidentiality standards (e.g., GDPR, HIPAA, and proprietary IP agreements).

Push-to-Talk (PTT) and Click-to-Speak models eliminate this vulnerability by enforcing physical and driver-level hardware isolation. When the user is not actively holding the PTT trigger or clicking the voice command button, the audio ingestion handle is closed. No audio samples are captured into RAM, no circular buffers are maintained, and no network packets are dispatched. By pairing hardware-isolated audio ingestion with dynamic mathematical normalization, PRIVACY68 demonstrates that enterprise-grade acoustic privacy can be achieved without sacrificing conversational speed or accuracy.

\section{Comparative Analysis of Existing Systems}
To contextualize the architectural contributions of PRIVACY68, Table \ref{tab:literature_comparison} presents an extensive comparative evaluation across six distinct categories of voice assistants and agentic automation platforms.

\begin{table}[htbp]
\centering
\scriptsize
\caption{Comprehensive Architectural Comparison of Existing Voice \& Agent Systems vs. PRIVACY68}
\label{tab:literature_comparison}
\begin{tabularx}{\textwidth}{|p{2.2cm}|X|X|X|X|p{2.6cm}|}
\hline
\textbf{Dimension} & \textbf{Commercial (Alexa/Siri)} & \textbf{Open Voice (Mycroft/OVOS)} & \textbf{AutoGPT / LangChain} & \textbf{Open-Interpreter} & \textbf{Proposed (PRIVACY68)} \\ \hline
\textbf{Microphone Isolation} & Always-On (Continuous Cloud Buffering) & Always-On (Local Wake-Word VAD) & N/A (Text-based Web / Terminal) & Manual Terminal Input & \textbf{100\% Click-to-Speak (Hardware Driver Isolated)} \\ \hline
\textbf{Acoustic Normalization} & Proprietary Cloud AGC & Energy Threshold Cutoff & None & None & \textbf{Dynamic Vector Peak Scaling ($\times 0.90$)} \\ \hline
\textbf{Primary ASR Engine} & Proprietary Cloud Acoustic Models & PocketSphinx / Vosk / Local Whisper & N/A (External Text API) & Local / Cloud Whisper & \textbf{Google Web Speech (Zero-Key) + Faster-Whisper Fallback} \\ \hline
\textbf{Reasoning Backend} & Closed Monolithic Cloud Intent Parser & Rule-based Regex \& Adapt Intent Parser & Single API Key (OpenAI / Anthropic) & Single LLM Provider (CLI configured) & \textbf{Multi-Provider Universal Core (Gemini, Groq, OpenRouter, Ollama)} \\ \hline
\textbf{Intra-Provider Failover} & None (Monolithic) & None (Fails on unhandled intent) & Manual Exception Throw & Manual Retries & \textbf{Autonomous Dynamic Failover Cascade} \\ \hline
\textbf{OS Automation Scope} & Proprietary Skills / Smart Home Only & Pre-defined Python Skills & Web Search / Shell Tools & Python Execution in Terminal & \textbf{Arbitrary Win32 / PowerShell / Python Scripting} \\ \hline
\textbf{Error Correction} & Generic ``I didn't understand'' & Skill Failure Exception & Re-prompting without OS sandbox & Manual User Prompting & \textbf{Autonomous 3-Pass Self-Healing Code Loop} \\ \hline
\textbf{Security Architecture} & Closed Cloud Sandbox & Ambient OS Privileges & Unrestricted CLI Execution & CLI Confirmation Prompt (\texttt{y/n}) & \textbf{Deterministic 3-Tier Security Guard + 15s Timeout Gate} \\ \hline
\textbf{Workflow Memory} & Proprietary Cloud History & Static Skill Repo & Vector Database (ChromaDB / Pinecone) & Ephemeral Chat Context & \textbf{Lightweight Deterministic Skill DB (\texttt{skills.json})} \\ \hline
\textbf{Presentation Interface} & Smart Speaker Hardware / Native Mobile UI & Dedicated Hardware / Linux GUI & Headless CLI / Web UI & Terminal Command Line & \textbf{Glassmorphic WebView2 HUD, Tray, \& LAN Mobile Remote} \\ \hline
\end{tabularx}
\end{table}

\section{Summary of Gaps Identified}
The synthesis of existing literature and empirical platforms reveals four major architectural and security gaps:
\begin{enumerate}
    \item \textbf{The Acoustic Privacy Deficit:} Existing voice assistants force an intolerable compromise: users must either accept perpetual microphone surveillance and background telemetry, or forgo voice-driven automation altogether.
    \item \textbf{The Brittleness of Handcrafted Grammars:} Conventional local voice platforms (such as Mycroft or native OS speech recognizers) depend on rigid regex patterns and slot-filling grammars, rendering them incapable of synthesizing novel operating system automation scripts.
    \item \textbf{Single-Vendor Vulnerability in AI Agents:} Emerging agentic frameworks are tightly coupled to a single proprietary cloud API, creating severe vulnerabilities to network outages, token exhaustion, and abrupt model deprecation.
    \item \textbf{Unchecked Security Risks in Agentic Subprocesses:} LLM-driven automation tools lack structured, policy-driven security architectures. They execute generated scripts with unrestricted user privileges, failing to protect the host workstation from destructive operations.
\end{enumerate}

PRIVACY68 systematically resolves each of these deficiencies through a unified, modular architecture detailed in the following chapters.
